\documentclass[10pt]{beamer}
\usetheme{Madrid}
\usepackage[T1]{fontenc}
\usepackage[utf8]{inputenc}
\usepackage[brazil]{babel}
\usepackage{lmodern}
\usepackage{hyperref}
\usepackage{listings}
\usepackage{xcolor}
\setbeamertemplate{navigation symbols}{}

\lstset{
  language=C++,
  basicstyle=\ttfamily\small,
  keywordstyle=\color{blue},
  commentstyle=\color{gray},
  stringstyle=\color{red},
  breaklines=true,
  showstringspaces=false,
}

\title{Comunicação Serial}
\subtitle{Curso de Microcontroladores}
\author{}
\date{}

\begin{document}

\begin{frame}
  \titlepage
\end{frame}

\begin{frame}{O que é comunicação serial?}
  \begin{itemize}
    \item Forma de enviar dados entre Arduino e computador.
    \item Usa 1 pino RX (recebe) e 1 pino TX (transmite).
    \item Velocidade padrão: 9600 baud.
    \item Essencial para debug e monitoramento.
  \end{itemize}
\end{frame}

\begin{frame}[fragile]{Inicializando comunicação}
\begin{lstlisting}
void setup() {
  Serial.begin(9600);  // Inicia serial a 9600 baud
  Serial.println("Sistema iniciado!");
}

void loop() {
  Serial.print("Valor: ");
  Serial.println(analogRead(A0));
  delay(1000);
}
\end{lstlisting}

  \begin{itemize}
    \item `Serial.begin()` inicia a comunicação.
    \item Deve ser chamado em `setup()`.
    \item Velocidade deve ser a mesma do monitor.
  \end{itemize}
\end{frame}

\begin{frame}{Funções Serial comuns}
  \begin{itemize}
    \item `Serial.print()`: imprime sem quebra de linha.
    \item `Serial.println()`: imprime com quebra de linha.
    \item `Serial.write()`: envia um byte bruto.
    \item `Serial.read()`: lê um byte do computador.
  \end{itemize}
\end{frame}

\begin{frame}[fragile]{Monitorando sensores}
\begin{lstlisting}
void loop() {
  int sensor1 = analogRead(A0);
  int sensor2 = analogRead(A1);
  int sensor3 = analogRead(A2);
  
  Serial.print("Sensores: ");
  Serial.print(sensor1);
  Serial.print(", ");
  Serial.print(sensor2);
  Serial.print(", ");
  Serial.println(sensor3);
  
  delay(100);
}
\end{lstlisting}

  Abra o Monitor Serial no Arduino IDE (Ctrl+Shift+M) para ver os valores.
\end{frame}

\begin{frame}[fragile]{Recebendo dados do computador}
\begin{lstlisting}
void loop() {
  if (Serial.available() > 0) {
    char comando = Serial.read();
    
    if (comando == 'L') {
      digitalWrite(LED_PIN, HIGH);
    } else if (comando == 'D') {
      digitalWrite(LED_PIN, LOW);
    }
  }
}
\end{lstlisting}

  \begin{itemize}
    \item `Serial.available()` verifica se há dados.
    \item `Serial.read()` pega um byte.
  \end{itemize}
\end{frame}

\begin{frame}{Debug em tempo real}
  \begin{itemize}
    \item Mostre valores de sensores regularmente.
    \item Indique quando estados mudam.
    \\item Facilita encontrar problemas rapidamente.
    \\item Não deixe muitos `println()` no código final (fica lento).
  \\end{itemize}
\\end{frame}

\\begin{frame}[fragile]{Dica: formatar saída}
\\begin{lstlisting}
void loop() {
  int valor = analogRead(A0);
  
  // Mostrar em diferentes formatos
  Serial.print("Dec: ");
  Serial.print(valor);
  Serial.print(", Hex: ");
  Serial.print(valor, HEX);
  Serial.print(", Bin: ");
  Serial.println(valor, BIN);
}
\\end{lstlisting}

  \\begin{itemize}
    \\item `DEC`: decimal (padrão)
    \\item `HEX`: hexadecimal
    \\item `BIN`: binário
  \\end{itemize}
\\end{frame}

\\begin{frame}{Importante para o seguidor}
  \\begin{itemize}
    \\item Monitore leitura dos sensores continuamente.
    \\item Veja valores do PID em tempo real.
    \\item Teste antes de rodar o robô na linha.
    \\item Serial é ferramenta de desenvolvimento, não coloque no código final otimizado.
  \\end{itemize}
\\end{frame}

\\end{document}
